\section{The entire spectral interval and a fixed strip}\label{sa:spectral}
The local spectral decrease uses the actual signed transfer, not
coefficientwise domination by a positive endpoint. From
Lemmas~\ref{lc:slow-transfer} and~\ref{lc:cofactor-max}, for every terminal row $N/2+O(1)$,
\begin{equation}\label{sa:real-upper}
 \limsup\frac1N\log|R_j(b,u(\delta))|
                  \le\max\{Q_b(\delta),\tfrac12\log b\}
                       \quad(0\le\delta\le1/2).
\end{equation}
For $b=0$ the second branch is omitted. The bound is uniform on the
whole interval. By \eqref{sa:Q-derivative}, $Q_b$ decreases strictly
when $b>0$ or $xy>0$. Also $Q_b(0)>\tfrac12\log b$ in that case:
the endpoint load is not the maximizer of the strictly concave entropy
when another entry is positive. The possible endpoint logarithmic
singularity at $x=y,\delta=1/2$ is integrable. Thus a separated
$|\delta|\ge\eta$ gives a fixed loss in every nonmonomial triangle.

We specify the small complex face extension used in sewing this bound.
In the proof of Lemma~\ref{lc:slow-transfer}, first remove a small initial
interval, then choose a fixed block length so the Cayley--Hamilton
bound for each two-dimensional transfer power is within any prescribed
exponential error of its spectral radius. Only after choosing that
length, perturb the face parameters. Matrix powers over this fixed
length vary continuously, so the same block bound persists. The
initial interval has the same integrable logarithmic allowance.
The scalar unused products are unchanged by face phases. Hence
\eqref{sa:real-upper}, with any prescribed error, holds on a fixed
sufficiently small complex face neighborhood. This argument neither
inverts an eigenbasis at a repeated root nor assumes phase positivity.

Fix the original radii $\rho$. On the small common-face-phase
neighborhood, one separated spectral coordinate supplies the above
loss in its nonmonomial triangle; the other factors have their
complete upper rates with errors whose finite sum is smaller than
half that loss. Equations \eqref{sa:even}--\eqref{sa:odd} retain both
triangle cofactors in each term, and their remaining prefactors are
bounded. On the complementary face torus use
\eqref{sa:global-angular}, uniform for every real spectral vector.
The Cauchy target factor is the same $\prod\rho_f^{-d}$ in both
regions. Its complete rate is precisely $\Phi$ from
\eqref{sa:stationarity}. We have therefore proved a real-cube gap
for the shared coefficient, not a sum of separately optimized factors.
The same argument without a loss, or the coefficientwise central
comparison, gives
\begin{equation}\label{sa:global-real}
 \sup_{\delta\in[-1/2,1/2]^R}|C_N(\delta)|
                                      \le e^{N(\Phi+o(1))}.
\end{equation}

For clarity a final elementary argument supplies a fixed complex strip
directly for this shared polynomial. A Laurent polynomial $F$ of
degree at most $d$ in each angular coordinate has, in either
half-plane of each variable, the subharmonic estimate
\begin{equation}\label{sa:Poisson}
 \log|F(x+\ii y)|\le
 (\mathcal P_y*\log|F|)(x)+2\pi d\sum_e|y_e|.
\end{equation}
Here $\mathcal P_y$ is the product of the circle Poisson kernels
at radii $e^{-2\pi|y_e|}$; a zero height denotes a point mass.
To prove it, multiply by the appropriate $z_e^d$ in each coordinate,
apply the disk subharmonic inequality successively, and subtract the
known monomial logarithms. Zeros follow by regularization, or by
using a continuous upper majorant on the boundary.

Apply the real gap with separation $\eta/2$ and
\eqref{sa:global-real}. If $|x_j|\ge\eta$, the one-coordinate
Poisson mass in $|t_j|\ge\eta/2$ tends uniformly to one as
$|y_j|\to0$. Choose the strip so it is at least $3/4$, and so
$2\pi R|y|$ and the $o(1)$ term together cost less than a quarter
of the real loss. Inequality \eqref{sa:Poisson} proves
\eqref{sa:H2}, possibly with a smaller constant. Cauchy's formula on
smaller disks and strips gives the fixed derivative assertions.
This completes Theorem~\ref{sa:main}.

